% =====================================================================
\chapter{Kiểm thử}\label{chap:kiemthu}
% =====================================================================

\section{Chiến lược kiểm thử}
\begin{figure}[H]
\centering
\mmd[0.92][0.6]{testpyramid}
\caption{Ba tầng kiểm thử: càng xuống dưới càng nhanh, tự động và chạy thường xuyên}\label{fig:testpyramid}
\end{figure}

\section{Kiểm thử tự động}
Bộ kiểm thử nằm trong \file{tests/} và chạy bằng \code{make -C tests}. Các module
nghiệp vụ của máy chủ được biên dịch với hai header giả (critical section và đồng hồ
điều khiển được). Hai sketch Arduino được biên dịch \emph{nguyên vẹn} cùng lớp giả lập
thư viện (\file{tests/arduino/}): đồng hồ \code{millis()}, mức logic từng chân, ADC,
Wi-Fi, UDP, đầu đọc MFRC522 và DHT11.

\begin{table}[H]
\centering\small
\caption{Các bộ kiểm thử tự động và kết quả}\label{tab:unittests}
\begin{tabularx}{\linewidth}{@{}lXr@{}}
\toprule
\textbf{Bộ} & \textbf{Nội dung kiểm tra} & \textbf{Kết quả} \\ \midrule
\file{json\_parsers} & Bản tin hợp lệ (mọi thứ tự khóa, khoảng trắng, giới hạn dải); 22 bản tin \code{parking\_update} sai và 16 bản tin \code{rfid\_event} sai đều bị từ chối (thiếu/thừa khóa, khóa lặp, sai kiểu, ngoài dải, ký tự lạ, dữ liệu thừa, chuỗi có ký tự thoát) & 57/57 \\
\file{rfid\_session} & Phí theo giờ (biên 0, 3600, 3601\,s, bão hòa); vào/ra, sai cổng; hóa đơn hết hạn đúng 60\,s; bãi đầy, tái sử dụng chỗ; nhiệt độ; nhật ký vòng (thứ tự, chỉ giữ 10); thông báo từ chối 5\,s & 74/74 \\
\file{gif\_frames} & 28 khung ảnh động giải nén đủ; nền đã xóa và byte đã đảo đúng; từ chối yêu cầu sai kích thước & 89/89 \\
\file{sensor\_client} & Chạy \file{SensorClient.ino}: nhắc lại đúng 1 gói/giây; gửi ngay sau 50\,ms khi có xe; lọc nhiễu 30\,ms; thứ tự S1–S5; báo cháy và hết cháy; quy đổi mực nước; mất/có lại Wi-Fi; mọi gói được bộ đọc của máy chủ chấp nhận & 15/15 \\
\file{rfid\_client} & Chạy \file{RFIDClient.ino}: UID hex in hoa; cổng vào/ra; chống lặp 2\,s trên cả hai đầu đọc; UID 10 byte; DHT11 lỗi giữ giá trị cũ; mọi gói được bộ đọc của máy chủ chấp nhận & 15/15 \\ \midrule
\textbf{Tổng} & & \textbf{250/250} \\
\bottomrule
\end{tabularx}
\end{table}

\begin{lstlisting}[language=C,caption={[Ví dụ kiểm thử chống rung của Sensor Client]Ví dụ kiểm thử chống rung của Sensor Client (\file{tests/test\_sensor\_client.cpp})}]
/* A car on S1 is reported as soon as the 50 ms debounce passes. */
fake::pin_level[SLOT_PINS[0]] = LOW;
run_for(40);
CHECK(fake::udp_sent.empty());            /* still inside the debounce window */
run_for(30);
CHECK(fake::udp_sent.size() == 1);
CHECK(last_packet(&s) && s.slots[0] == 1); /* parsed by the real server parser */
\end{lstlisting}

\begin{lydo}[Tại sao kiểm thử trên máy tính]
Lỗi trong bộ đọc bản tin, cách tính phí hay logic chống rung rất khó tái hiện trên
phần cứng (phải đúng thời điểm, đúng chuỗi byte). Trên máy tính, đồng hồ giả cho phép
kiểm tra chính xác đến từng mili giây, cả bộ chạy trong vài giây và chạy lại sau mỗi
lần sửa mã. Việc đưa gói tin của client vào bộ đọc thật của máy chủ bảo đảm hai bên
luôn khớp định dạng — lỗi kiểu ``client gửi \code{temperature\_c} trong
\code{parking\_update} bị máy chủ từ chối'' được phát hiện ngay.
\end{lydo}

Bộ kiểm thử cũng đã giúp phát hiện một lỗi \emph{trong chính bài kiểm thử}: khi kiểm
tra ``bãi đầy'', một thẻ từ bài kiểm thử trước vẫn còn trong bãi khiến chỉ còn 4 chỗ
trống. Bài kiểm thử được sửa để tự dọn trạng thái — đúng như mục đích của việc chạy
kiểm thử thường xuyên.

\section{Kiểm thử tích hợp bằng công cụ giả lập}
Các lệnh ở Bước~10 (Chương~\ref{chap:huongdanchay}) gửi bản tin thật tới máy chủ thật;
dùng để kiểm tra màn hình, thanh chắn, còi, trang web và nhật ký khi chưa có đủ client.

\section{Kịch bản kiểm thử hệ thống trên phần cứng}
\begin{longtable}{@{}p{0.9cm}>{\raggedright\arraybackslash}p{5.6cm}>{\raggedright\arraybackslash}p{7.9cm}@{}}
\caption{Kịch bản kiểm thử trên phần cứng và hành vi mong đợi}\label{tab:systest}\\
\toprule \textbf{\#} & \textbf{Thao tác} & \textbf{Hành vi mong đợi} \\ \midrule \endfirsthead
\toprule \textbf{\#} & \textbf{Thao tác} & \textbf{Hành vi mong đợi} \\ \midrule \endhead
S1 & Cấp nguồn RFID Client, rồi máy chủ & Có Wi-Fi \code{ESP32-Parking}; máy chủ báo \code{Got IP} và \code{mDNS: parking.local -> ...}; TFT hiện giao diện; hai thanh chắn đóng \\
S2 & Cấp nguồn Sensor Client, mở \code{parking.local} & Web hiện sơ đồ 5 ô, \code{LIVE / SENSOR} \\
S3 & Đặt xe vào ô S2 & TFT đổi ô S2 sang đỏ gần như tức thì; web hiện hình xe ở S2 trong $\le$ 1\,s \\
S4 & Rê vật nhanh qua cảm biến ($<$ 50\,ms) & Trạng thái không đổi \\
S5 & Quẹt thẻ mới ở cổng vào & Thanh chắn vào mở, tự đóng sau 3\,s; web có dòng \code{IN}; nhiệt độ cập nhật \\
S6 & Giữ thẻ trên đầu đọc 1\,s & Chỉ một sự kiện được gửi \\
S7 & Quẹt lại thẻ đó ở cổng vào & Không mở; TFT \code{WRONG GATE – ALREADY INSIDE - USE EXIT GATE} 5\,s; web \code{REJECTED · WRONG GATE} \\
S8 & Quẹt thẻ đó ở cổng ra & Thanh chắn ra mở; TFT hiện hóa đơn 60\,s; web \code{OUT} kèm thời lượng và phí \\
S9 & Đưa lửa gần cảm biến lửa & Còi kêu; TFT \code{DANGER}; hai thanh chắn mở; băng rôn đỏ trên web \\
S10 & Bấm \code{CLOSE} trên web khi đang cháy & Bị từ chối, thanh chắn vẫn mở \\
S11 & Tắt lửa & Còi tắt, hai thanh chắn đóng, TFT trở lại bình thường \\
S12 & Bấm \code{OPEN}/\code{CLOSE}/\code{EMERGENCY} trên web & Thanh chắn tương ứng chuyển động, trạng thái trên web đúng \\
S13 & Quẹt thẻ thứ 6 khi đã có 5 xe & Không mở; TFT \code{PARKING FULL} \\
S14 & Tắt RFID Client rồi bật lại & Máy chủ và Sensor Client tự kết nối lại; \code{parking.local} trỏ đúng IP mới \\
S15 & Gửi gói UDP rác từ máy tính & Máy chủ ghi cảnh báo và tiếp tục hoạt động \\
\bottomrule
\end{longtable}
